\section{Linked Study Packages}
The public PDF is the research-program overview, not the complete protocol archive for every bounded claim. Study-specific repositories carry the exact preprocessing boundary, frozen split design, result tables, and reproduction scripts. I use logical method names in the text while preserving legacy repository slugs where those slugs are the working public URLs.

\begin{table*}[t]
\centering
\caption{Focused study packages linked from the main research PDF.}
\small
\begin{tabular}{p{0.27\linewidth}p{0.46\linewidth}p{0.21\linewidth}}
\toprule
Package & What it contains & Link \\
\midrule
Paired-Acquisition Neural Factorization external canine SCC validation & Independent five-scanner canine SCC paired-acquisition validation on 805 geometry-qualified five-view regions from 44 biological samples. The locked projection reduced sample-blocked scanner-probe accuracy by 0.380 absolute while preserving same-region retrieval. & \href{https://github.com/matthewvaishnav/pathoalign-external-caninescc}{Repository} \\
Paired-Acquisition Neural Factorization allocation study & Resource-allocation study testing biological pair diversity versus anchor repetition under matched pair-presentation budgets. The strongest allocation contrast was 200 versus 50 pairs at budget 6,400, with positive budget-doubling effects across allocations. & \href{https://github.com/matthewvaishnav/pathoalign-pair-repeat-allocation}{Repository} \\
\bottomrule
\end{tabular}
\label{tab:linked_study_packages}
\end{table*}

Do not link to study-package GitHub Pages PDFs until those pages are actually published. The repository link is the stable public entry point for each study package.
